\chapter{Facial Droop}
\index{Facial Droop}

\section*{Definition and Overview}
Facial droop is the weakness or paralysis of the muscles on one or both sides of the face, causing the face to appear uneven, with difficulty closing the eye and smiling. It has two main causes: Bell's palsy, a temporary weakness of the facial nerve that usually recovers, and stroke, a medical emergency that requires immediate treatment. Because the two can look similar, any sudden facial droop should be treated as a possible stroke until proven otherwise. This chapter explains the causes, the differences, and what to do.

\section*{Epidemiology}
Facial droop is a common reason for urgent medical assessment. Stroke is a leading cause of death and disability in many countries, and facial droop is one of its classic signs. Bell's palsy affects about 1 in 5,000 people each year, most often adults, and is the most common cause of sudden facial weakness. Other causes, including infections and tumours affecting the facial nerve, are rarer.

\section*{Aetiology and Pathophysiology}
Facial droop arises from problems with the facial nerve or the brain.
\begin{itemize}
    \item \textbf{Bell's palsy:} inflammation of the facial nerve, possibly viral, causing one-sided facial weakness; the cause is not fully understood
    \item \textbf{Stroke:} interruption of the blood supply to the brain, damaging the areas that control facial movement
    \item \textbf{Other causes:} infections such as shingles, middle ear infections, Lyme disease, and rarely tumours or trauma
    \item \textbf{The facial nerve:} controls the muscles of expression on each side of the face
    \item \textbf{Differences:} in stroke, the weakness usually spares the forehead, and other symptoms such as arm weakness and speech problems are present; in Bell's palsy, the whole side of the face is weak, including the forehead
    \item \textbf{Urgency:} stroke treatment works only if given quickly, so the priority is to rule out stroke
\end{itemize}

\section*{Clinical Features}
The signs depend on the cause.
\begin{itemize}
    \item One side of the face droops, with difficulty smiling
    \item Difficulty closing the eye on the affected side
    \item Drooling and difficulty eating and drinking
    \item In Bell's palsy: sudden onset, pain behind the ear, a change in taste, and sensitivity to sound
    \item In stroke: sudden weakness, with arm weakness, speech difficulty, confusion, or vision loss
    \item In stroke: other symptoms may include dizziness, a severe headache, and difficulty walking
\end{itemize}

\section*{Differential Diagnosis}
The urgent distinction is between stroke and other causes.
\begin{itemize}
    \item Stroke and transient ischaemic attack, which are emergencies
    \item Bell's palsy, the common benign cause
    \item Ramsay Hunt syndrome, shingles affecting the facial nerve
    \item Middle ear infection causing facial weakness
    \item Tumours and trauma affecting the nerve
    \item Rarely, other neurological conditions
\end{itemize}

\section*{When to Seek Medical Care}
Sudden facial droop requires immediate assessment, because it may be a stroke. Call for an ambulance if there are any stroke symptoms, and seek urgent review even if the droop seems mild.

\textbf{Contact your local emergency services for:}
\begin{itemize}
    \item Sudden facial drooping, especially with arm weakness or difficulty speaking
    \item Sudden confusion, vision loss, or severe headache
    \item Sudden weakness or numbness of the limbs
    \item Difficulty walking or dizziness with facial droop
\end{itemize}

\section*{Diagnosis and Investigations}
Diagnosis distinguishes stroke from Bell's palsy.
\begin{itemize}
    \item \textbf{Clinical assessment:} the pattern of weakness and the presence of other symptoms
    \item \textbf{Stroke assessment:} examination of the face, arms, and speech
    \item \textbf{Imaging:} CT or MRI of the brain for suspected stroke
    \item \textbf{Blood pressure and tests:} as part of the stroke assessment
    \item \textbf{Bell's palsy:} a clinical diagnosis made after stroke is excluded
\end{itemize}

\section*{Management}
Management depends on the cause.
\begin{itemize}
    \item \textbf{Stroke:} urgent treatment in hospital, including clot-dissolving treatment where suitable, followed by rehabilitation
    \item \textbf{Bell's palsy:} corticosteroids, started early, which improve recovery; eye protection and drops while the eye cannot close
    \item \textbf{Other causes:} treatment of the infection or underlying condition
    \item \textbf{Supportive care:} protection of the eye, gentle facial exercises, and pain relief
    \item \textbf{Follow-up:} review of recovery
\end{itemize}

\section*{Complications}
\begin{itemize}
    \item In stroke: permanent weakness, disability, and further events without treatment
    \item In Bell's palsy: incomplete recovery, a dry or damaged eye, and rarely permanent weakness
    \item Difficulty eating, drinking, and speaking
    \item The psychological impact of a changed appearance
\end{itemize}

\section*{Prognosis and Outcomes}
The outlook depends on the cause. Stroke outcomes improve dramatically with rapid treatment, so every minute counts. Bell's palsy recovers fully in the majority of cases, usually within weeks to months, and corticosteroids improve the odds of complete recovery. With appropriate care, most people regain facial function.

\section*{Prevention and Self-Care}
\begin{itemize}
    \item Know the signs of stroke and act F.A.S.T.: Face, Arms, Speech, Time
    \item Contact your local emergency services immediately for any suspected stroke
    \item Manage blood pressure, cholesterol, diabetes, and other stroke risk factors
    \item In Bell's palsy: protect the eye with drops and a patch at night
    \item Eat soft foods and chew on the unaffected side if needed
    \item Attend follow-up and rehabilitation
    \item Practise gentle facial exercises as advised
\end{itemize}

\section*{Sources and Further Reading}
The following reputable sources provide further information for healthcare students and patients seeking reliable, current advice about facial droop.
\begin{itemize}
    \item Stroke Foundation Australia. \textit{F.A.S.T. signs of stroke.} https://strokefoundation.org.au/
    \item Healthdirect Australia. \textit{Bell's palsy.} https://www.healthdirect.gov.au/bells-palsy
    \item NHS. \textit{Bell's palsy.} https://www.nhs.uk/conditions/bells-palsy/
    \item Mayo Clinic. \textit{Facial paralysis and Bell's palsy.} https://www.mayoclinic.org/diseases-conditions/bells-palsy/
    \item MedlinePlus. \textit{Facial nerve palsy.} https://medlineplus.gov/ency/article/003028.htm
    \item MSD Manual. \textit{Bell's palsy.} https://www.msdmanuals.com/
\end{itemize}
